\subsection{\texorpdfstring{Commutative rings of characteristic \(p\)}{Commutative rings of characteristic p}}

In this No.~and the following one p denotes a prime number.

THEOREM 2. --- Let A be a commutative ring of characteristic p.~The mapping \(a \mapsto a^{p}\) is an endomorphism of the ring A, that is, we have the relations

\[
(a + b) ^ {p} = a ^ {p} + b ^ {p}\tag{1}
\]

\[
(a b) ^ {p} = a ^ {p} b ^ {p}\tag{2}
\]

for \(a, b\) in A.

Formula (2) follows from the commutativity of A. To prove (1) we use the binomial formula \((a + b)^p = a^p + b^p + \sum_{i=1}^{p-1} \binom{p}{i} \cdot a^ib^{p-i}\) ; since \(p \cdot x = 0\) for all \(x \in \mathbf{A}\) , it suffices to prove the following lemma:

Lemma 1. --- Let \(p\) be a prime number and \(i\) an integer in the range from 1 to \(p - 1\) , then the binomial coefficient \(\binom{p}{i}\) is an integer divisible by \(p\) .

We argue by induction on \(i\) , the case \(i = 1\) being immediate from the formula \(\binom{p}{1} = p\) . Suppose that \(2 \leqslant i \leqslant p - 1\) and that \(\binom{p}{i - 1}\) is divisible by \(p\) . Then the integer \(i\binom{p}{i} = (p - i + 1)\binom{p}{i - 1}\) belongs to the prime ideal \(p\mathbf{Z}\) of \(\mathbf{Z}\) ; since \(i \notin p\mathbf{Z}\) , we have \(\binom{p}{i} \in p\mathbf{Z}\) and the lemma follows.

Let A be a commutative ring of characteristic p and f an integer \(\geqslant0\) . From Th. 2 we deduce by induction on f that the mapping \(a\mapsto a^{p^{f}}\) is an endomorphism of the ring A. In particular we have the relation

\[
(a _ {1} + \dots + a _ {n}) ^ {p ^ {f}} = a _ {1} ^ {p ^ {f}} + \dots + a _ {n} ^ {p ^ {f}}\tag{3}
\]

for any \(a_{1},\ldots,a_{n}\) in A. The mapping \(a\mapsto a^{p}\) is sometimes called the Frobenius endomorphism of A. Taking \(\mathbf{A}=\mathbf{F}_{p}\) and \(a_{i}=1\) , we obtain from (3) the relation:

\[
n ^ {p ^ {f}} \equiv n \mod p (n \in \mathbf {Z}, f \in \mathbf {N}).\tag{4}
\]

For each subset S of A we denote by \(S^{p^{f}}\) the set of elements of A of the form \(x^{p^{f}}\) with \(x \in S\) \(^{1}\) . In particular, if K is a subring of A, the set \(K^{p^{f}}\) is a subring of A. If K is a subring of A and S a subset of A we denote by K{[}S{]} the subring of A generated by \(K \cup S\) ; when A is a field, we denote by \(\mathrm{K}(S)\) the field of fractions of K{[}S{]}, that is, the subfield of A generated by \(K \cup S\) .

PROPOSITION 2. --- Let A be a commutative ring of characteristic p, K a subring of A, S a subset of A and f a positive integer.

\begin{enumerate}
\def\labelenumi{\alph{enumi})}
\item
  We have \(\mathbf{K}[\mathbf{S}]^{p^f} = \mathbf{K}^{p^f}[\mathbf{S}^{p^f}]\) , and if \(\mathbf{A}\) is a field, \(\mathbf{K}(\mathbf{S})^{p^f} = \mathbf{K}^{p^f}(\mathbf{S}^{p^f})\) .
\item
  If the K-module K{[}S{]} is generated by the family \((a_{i})_{i\in I}\) of elements of A, then the K-module \(K[S^{p^{f}}]\) is generated by the family \((a_{i}^{p^{f}})_{i\in I}\) .
\end{enumerate}

Since K{[}S{]} is the subring of A generated by K ∪ S, its image K{[}S{]} \(^{p^{f}}\) under the endomorphism \(\pi: a \mapsto a^{p^{f}}\) of the ring A is the subring of A generated by the image K \(^{p^{f}}\) ∪ S \(^{p^{f}}\) of K ∪ S under \(\pi\) , whence K{[}S{]} \(^{p^{f}}\) = K \(^{p^{f}}\) {[}S \(^{p^{f}}\) {]}. The case of fields is treated similarly; this proves a).

It is clear that the family \((a_i^{p^f})_{i\in I}\) generates the \(\mathbf{K}^{p^f}\) -module \(\mathbf{K}[\mathbf{S}]^{p^f}\) . The K-module \(\mathbf{K}[\mathbf{S}^{p^f}]\) is generated by products of the form \(x_1^{p^f}\ldots x_n^{p^f} = (x_1\ldots x_n)^{p^f}\) with \(x_1,\ldots ,x_n\) arbitrary in S, hence also by the set \(\mathbf{K}[\mathbf{S}]^{p^f}\) . Assertion b) follows directly from this.

\(^{1}\) Of course the set \(S^{p^{f}}\) should not be confused with the set product of \(p^{f}\) sets equal to S, nor with the set of products of \(p^{f}\) elements belonging to S.

